\documentclass[a4paper]{article}
 
\usepackage{polski}
\usepackage[utf8]{inputenc} 
\usepackage[pdftex]{graphicx}
\usepackage[left=2.5cm,top=2.5cm,right=2.5cm,bottom=2.5cm,nohead,nofoot]{geometry}
\usepackage{fancyhdr}
\usepackage{indentfirst} 

\renewcommand{\maketitle}{
\begin{titlepage}
\begin{center}
\huge
\textbf{Akademia Górniczo-Hutnicza}\\
\Large
Wydzial Elektrotechniki, Automatyki, Informatyki i Elektroniki\\
\vspace*{0.4cm}
Katedra Informatyki
\vspace*{0.2cm}
\end{center}
\begin{figure}[!ht]
\centering
\includegraphics[scale=0.5]{Grafika/agh-logo.png}
\end{figure}
\vspace*{0.1cm}
\begin{center}
\LARGE
\textbf {Implementacja algorytmu karalucha}\\
zastosowanie algorytmu do rozwiązywania zagadnień optymalizacji kombinatorycznej\\
\vspace*{0.8cm}
\begin{table}[!htp]
\hspace*{0.45cm}
\begin{tabular}{p{3cm}p{9cm}}
\textit{Autorzy:}& \parbox{4cm} {\textbf {Łukasz Adamczyk\\Piotr Adamczyk\\Bartosz Błażewicz\\Mariusz Cybula\\Piotr Dereszowski\\Dawid Marszałek\\Marcin Piechota\\Jakub Warzocha\\Joanna Wojtas\\Rafał Szczodruch\\Kamil Nowak}}\\
\\
\textit{Kierunek studiów:} &\textbf {Informatyka, II rok} \\
\end{tabular}
\end{table}
\vspace{0.5cm}
\center{\small Kraków 2012}
\end{center}
\end{titlepage}
}

\begin{document}
\maketitle
\tableofcontents
\newpage

\section{Informacje ogólne}

\subsection{Wprowadzenie - Algorytmy inspirowane naturą}
Algorytmy inspirowane naturą w ostatnich latach zyskały dużą popularność ze względu na fakt,
że wiele rzeczywistych problemów optymalizacyjnych stały się bardziej dynamiczne oraz złożone.
Rozmiar i złożoność problemów w dzisiejszych czasach wymaga rozwoju wielu metod,
których efektywność jest mierzona na podstawie zdolności do znalezienia zadowalających wyników w rozsądnym czasie,
a nie tylko do zagwarantowania optymalnego rozwiązania. Celem naukowców stała się więc obserwacja natury,
próba zrozumienia efektywnych algorytmów wykorzystywanych przez inne organizmy żywe.
Przeniesienie tej wiedzy i dopasowanie rozwiązań na grunt informatyki - specjaliści oczekują uzyskania jak najlepszych rezultatów.

\subsection{Karaluch - co o nim wiemy}
\subsubsection{Co to jest karaluch?}
%\begin{table}[!htp]
%\begin{tabular}{p{3cm}p{9cm}}
%{Karaluch:}&{karaczan wschodni}\\
%{Gromada:}&{owady}\\
%{Podgromada:}&{uskrzydlone}\\
%{Rząd:}&{karaczany}\\
%\end{tabular}
%\end{table}
Niewielki owad żyjący w siedliskach człowieka. Żyje w kuchniach, magazynach, wyjadając produkty żywnościowe. Długość ciała dochodzi do 3 cm. Ubarwienie czerwone do czarno-brązowego.
Ciało owadów pokryte jest cienką, woskową powłoką, nie przepuszczającą wody. Czułki są zazwyczaj tej samej długości co ciało i stanowią główny narząd zmysłu. Karaczany mają narządy gębowe typu gryzącego. Na odwłoku występują często gruczoły wydzielające odstraszającą woń i ciecz. Pierwsza para skrzydeł jest skórzasta i tworzy pokrywę, druga para jest błoniasta.
Wykazują cechy stadne. Rozpoznają się wzajemnie dzięki specyficznemu wydzielanemu zapachowi. Podążają chaotycznie w miejsca gdzie są już inne osobniki. Niektóre gatunki prowadzą nocny tryb życia. Najchętniej gromadzą się w szparach i szczelinach oraz w miejscach wilgotnych. Mogą przenosić choroby, statystycznie rzecz biorąc, na każdym karaczanie przebywa nawet 80 chorobotwórczych bakterii.
\subsubsection{Ciekawostki} 
Karaluchy są wszystkożerne - są jednymi z niewielu zwierząt zdolnymi do strawienia keratyny - białka włosa.
W niektórych krajach ma nie tylko negatywny wizerunek w Rosji np. uznawany jest za dowód dostatku w domu, bo tam gdzie marnuje się jedzenie, tam są i karaluchy.
Niewymagający i niewybredny owad słynie z odporności, np. bez jedzenia może się obejść nawet 40 dni. Są one odporne na wiele trucizn i potrafią bez uszczerbku znieść promieniowanie radioaktywne w dawkach nawet 100 krotnie większych niż dawka szkodliwa dla człowieka. Znane są przypadki karaluchów żyjących mimo braku głowy
Zespół naukowców z Nottingham University odkrył w mózgach karaluchów i szarańczy aż 9 różnych substancji o silnym działaniu przeciwbakteryjnym, a siedem z nich zabijało 90 procent „superbakterii” , nie szkodząc przy tym ludzkim komórkom.

Z racji swojej odrazy do karaluchów ludzie nazywali znienawidzone owady tak jak nie lubiane nacje, np. jeden z krajowych gatunków nazywany jest w Polsce prusakiem (uznana nazwa naukowa!) lub popularnie szwabem.





\subsection{Problem Harmonogramowania}
Obecnie jednym z podstawowych elementów zarządzania wszelkiego typu projektami jest proces harmonogramowania. Pod
"wszelkiego typu" możemy rozumieć niemalże każde przedśiewzięcie - inżyniera, sektor IT, inwestycje, przemysł transportowy
, budowlany itp. Za terminem harmonogramownia stoi nic innego jak planowanie najbardziej optymalnej sekwencji operacji
oraz zminimalizownie czasu wykonania tychże sekwencji. Określamy skończone ramy czasowe wykonania danych zadań oraz
znajdujemy możliwe równoległe scieżki wykonywania operacji. Oczywistą korzyścią harmonogramowania jest najlepsze
wykorzystaniu czasu oraz możliwośc dokładnego obliczenia terminu wykonania całego lub części przedsięwzięcia.
W ostatnich latach powstało wiele pokaźnej literatury poruszającej problem, zrodziło się wiele definicji, modeli oraz teorii. 

\subsection{Problem komiwojazera}
Problem komiwojażera (bądź jak ktoś woli TSP - travelling salesman problem) jest jednym z najstarszych - a przynajmniej
jednym z najbardziej popularnych problemów optymalizacyjnych. Zapewne wiele osób z tym problemem już się spotkało i wie,
że jest związany z teorią grafów. Ściślej rzecz ujmując sprowadza się do znalezienia cyklu Hamiltona w grafie pełnym
(graf musi być pełny, aby istniało rozwiązanie).\\

\noindent Kilka słów wyjaśnień:

\begin{description}
\item [Graf pełny] - jest to graf w którym dla każdej pary wierzchołków istnieje krawędź je łącząca.
             W przypadku grafu pełnego o n wierzchołkach wiadomo, iż posiada dokładnie n(n-1)/2 krawędzi.
\item [Cykl Hamiltona] - jest to taki cykl, który zawiera każdy wierzchołek danego grafu dokładnie 1 raz (prócz pierwszego
                 wierzchołka, ponieważ w nim rozpoczyna się i kończy cykl).
\end{description}
Dodatkowo wiadomo, krawędziom grafu można nadawać różne wagi.\\

Wracając do problemu: dane jest n miast oraz odległość/koszt podróży/czas podróży między nimi. Nasz strudzony
komiwojażer musi odwiedzić wszystkie te miasta w celu dostarczenia przesyłek (cokolwiek stanowią) jednak z zachowaniem
zminimalizowanych kosztów. W zależności od problemu celem może być znalezienie najkrótszej/najtańszej/najszybszej trasy.
Warunek jest taki, że każdą miejscowość ma odwiedzić dokładnie 1 raz po czym ma wrócić do miejscowości w której "wojaże"
rozpoczął.\\

\noindent W odniesieniu do teorii grafów:
\begin{itemize}
\item wierzchołek grafu reprezentuje jedno miasto;
\item  krawędź reprezentuje drogę łączącą miasta;
\item  jako, że istnieją pewne kryteria (odległość/koszt/czas), krawędzie grafu mają przyporządkowane odpowiednie wagi.
\end{itemize}









\end{document}